% Part: sets-functions-relations
% Chapter: arithmetization
% Section: reflections
% Hindi translation (hi-Deva-IN) of Open Logic commit 9620cc73f9c8e0ad003c514a5d3748f29611c4c0.

\documentclass[../../../include/open-logic-section]{subfiles}

\begin{document}

\olfileid[hi]{sfr}{arith}{ref}
\olsection{कुछ दार्शनिक विचार}

तकनीकी विवरण तो इतने हुए। लेकिन उनसे मिला क्या?

एक लगभग निर्विवाद उपलब्धि सुंदर शुद्ध गणित है। इसके अतिरिक्त कुछ गहरी
वैचारिक उपलब्धियाँ भी हुईं। यह समझना गंभीर अंतर्दृष्टि थी कि पूर्णता गुणधर्म
वास्तविक और परिमेय संख्याओं के बीच निर्णायक अंतर व्यक्त करता है। फिर,
डेडेकिंड विच्छेदों के रूप में वास्तविक संख्याओं का स्पष्ट निर्माण गणितीय
विश्लेषण को ठोस आधार देता है। हम जानते हैं कि \emph{पूर्ण क्रमित क्षेत्र} की
संकल्पना सुसंगत है, क्योंकि विच्छेद ठीक ऐसा ही क्षेत्र बनाते हैं।

इतना सब होने पर भी इन उपलब्धियों के बारे में कुछ संशय सामने रखने चाहिए।

पहला, यह स्पष्ट नहीं है कि विच्छेदों के रूप में वास्तविक संख्याओं को सोचना,
उनके परिचित (संभवतः अनंत) दशमलव प्रसारों के रूप में सोचने की अपेक्षा
\emph{अधिक} कठोर है। वास्तविक संख्याओं का यह दूसरा ``निर्माण'' कौशी अनुक्रमों
के माध्यम से उनके निर्माण से कुछ मिलता-जुलता है; किंतु वास्तव में गणितज्ञ इसे
मूलतः सत्रहवीं शताब्दी के आरंभ से जानते थे (\olref[cauchy]{sec} देखें)।
कठोरता में वास्तविक वृद्धि इस बोध से आई कि वास्तविक संख्याओं में पूर्णता
गुणधर्म होता है; विशिष्ट समुच्चयों के रूप में वास्तविक संख्याओं को बनाने की
क्षमता अपने-आप में शायद उतनी रोचक नहीं है।

यह बात और भी कम स्पष्ट है कि पूर्णांकों या परिमेय संख्याओं का (कहीं अधिक सरल)
अंकगणितीकरण उन क्षेत्रों की कठोरता बढ़ाता है। यहाँ एक सरल बात ध्यान देने योग्य
है। प्राकृत संख्याओं के क्रमित युग्मों के तुल्यता वर्गों के रूप में पूर्णांकों
का \emph{निर्माण} करने, फिर पूर्णांकों के क्रमित युग्मों के तुल्यता वर्गों के
रूप में परिमेय संख्याओं का निर्माण करने, और फिर परिमेय संख्याओं के समुच्चयों
के रूप में वास्तविक संख्याओं का निर्माण करने के तुरंत बाद हम इन निर्माणों को
\emph{भूल जाते हैं}। विशेषतः किसी गणितीय प्रमाण के दौरान कोई भी इन निर्माणों
का \emph{उपयोग} नहीं करना चाहेगा (उस प्रमाण को छोड़कर जो यह दिखाए कि निर्माण
अपेक्षित ढंग से व्यवहार करते हैं)। किसी वास्तविक संख्या के बारे में सीधे
बोलना, प्राकृत संख्याओं के समुच्चयों के समुच्चयों के समुच्चयों के समुच्चयों के
समुच्चयों के समुच्चयों के समुच्चय के बारे में बोलने से बहुत आसान है।
%(For reals were constructed as sets of rationals; and these were constructed as equivalence classes of ordered pairs of integers, i.e., sets of sets of sets of integers; etc.).
%2/3 = [2,3]
%	i.e. {<2,3>, ... }
%	{ {{2},{2,3}}, ... }
%
%reals are
%	sets of rationals
%
%rationals are
%	sets of sets of sets of integers
%
%integers are
%	sets of sets of sets of naturals
	
सबसे अधिक संदेह इस बात पर है कि ये परिभाषाएँ हमें बताती हैं कि पूर्णांक,
परिमेय संख्याएँ या वास्तविक संख्याएँ \emph{हैं} क्या---\emph{तत्त्वमीमांसा की
दृष्टि से}। अर्थात यह संदेहास्पद है कि वास्तविक संख्याएँ (मसलन) कुछ विशेष
समुच्चय (समुच्चयों के समुच्चयों के समुच्चय\ldots) ही \emph{हैं}। ऐसे विचार
की मुख्य बाधा यह है कि निर्माण कई अलग-अलग ढंग से किया जा सकता था। वास्तविक
संख्याओं के मामले में कुछ सचमुच रोचक और भिन्न निर्माण हैं
(\olref[cauchy]{sec} देखें)। लेकिन अलग निर्माण पाने का एक बिल्कुल मामूली उपाय
यह है: \olref[sfr][rel][ref]{sec} की तरह हम क्रमित युग्म को थोड़ा अलग ढंग से
परिभाषित कर सकते थे; यदि क्रमित युग्म की इस वैकल्पिक संकल्पना का उपयोग किया
होता, तो हमारे निर्माण उतनी ही अच्छी तरह काम करते, पर अंत में हमें अलग वस्तुएँ
मिलतीं। इसलिए पूर्णांकों, परिमेय संख्याओं और वास्तविक संख्याओं के कई परस्पर
प्रतिस्पर्धी समुच्चय-सैद्धांतिक निर्माण हैं। अब यह दावा करना मनमाना (और
लज्जाजनक) होगा कि पूर्णांक (मसलन) \emph{इन} समुच्चयों को कहते हैं, न कि
\emph{उन} समुच्चयों को। (\olref[sfr][rel][ref]{sec} की तरह यह उस तर्क का एक
उदाहरण है जिसे \citealt{Benacerraf1965} ने प्रसिद्ध बनाया।)

एक और बात उठाना उचित है: हमारे निर्माणों में कुछ बहुत \emph{विचित्र} है। हमने
प्राकृत संख्याओं से आरंभ किया। फिर पूर्णांक बनाए, और ``पूर्णांकों का $0$'' बनाया,
अर्थात $ \equivrep{0,0}{\Intequiv}$। किंतु $0 \neq
\equivrep{0,0}{\Intequiv}$। वास्तव में, हमारे निर्माणों के अनुसार \emph{कोई
भी} प्राकृत संख्या पूर्णांक नहीं है। लेकिन यह सहज-बोध के अत्यंत प्रतिकूल लगता
है। वास्तव में \olref[sfr][set][imp]{sec} में हमने बिना अधिक तर्क के दावा किया
था कि $\Nat \subseteq \Rat$। यदि निर्माण हमें ठीक-ठीक बताते हैं कि संख्याएँ
\emph{हैं क्या}, तो वह दावा सीधे-सीधे असत्य था।

तो कुछ दूरी से देखने पर हम कहाँ पहुँचते हैं? सरल समुच्चय-सिद्धांत में काम करते
हुए और प्राकृत संख्याओं को उपलब्ध मानकर, हम पूर्णांकों, परिमेय संख्याओं और
वास्तविक संख्याओं को कुछ समुच्चयों के रूप में \emph{बरत} सकते हैं। इस अर्थ में
हम इन वस्तुओं के सिद्धांतों को समुच्चय-सिद्धांत में \emph{अंतःस्थापित} कर सकते
हैं। किंतु इस अंतःस्थापन का दार्शनिक अर्थ बिलकुल सीधा नहीं है।

निश्चय ही यह सब अंतिम निष्कर्ष नहीं है!{} बात केवल इतनी है। यह दिखाने के लिए
कि वास्तविक संख्याओं का अंकगणितीकरण गहरा दार्शनिक महत्त्व \emph{रखता है}, कुछ
अतिरिक्त \emph{दार्शनिक} तर्क चाहिए।

%Additionally: for the entire duration of this chapter, we have treated
%the natural numbers as simply \emph{given} to us. But what can we do
%with them? That is the subject matter for the next chapter.

\end{document}
